% Standalone input component; author draft, not an acceptance record.
% Base: V18 ca1ea465... or pending 800e23b0... .
% Sources and exact artifacts are bound in SOURCE-MANIFEST.json.
% This file does not contain or assert a Busy Beaver threshold theorem.
\section{The 198-coordinate family and its inherited inputs}\label{q198:main}

We use the unbalanced 198-coordinate table of
\cite{q198-Conant198}. Its input is $X$, and its remaining coordinates
are $Y_1,\ldots,Y_{197}$. Put
\[
 A_N=149^{2N+1},\qquad K_0=31581660^{28},\qquad
 H_0=(5898050446009651968K_0)^2\qquad(N\ge1).
\]
The parameter $A_N$ is a fixed integer coefficient, not an existential
coordinate. Let $r_{1,N},\ldots,r_{184,N}$ be the primitive integral
quadratic residuals of this table and put
\[
 P_N=\sum_{j=1}^{184}r_{j,N}^2\in
 \Z[X,Y_1,\ldots,Y_{197}].
\]
The table is the file \nolinkurl{circuits/inverse198/quadratic.json}
in the pinned Conant package; its expanded sum of squares is
\nolinkurl{circuits/inverse198/polynomial.json}.
The exact interpretation and finite coefficient identities are proved
in Section~\ref{q198:finite-certificate} below. All coordinate counts
include the input.

The proof uses the results already established in this manuscript
and listed in Section~\ref{q198:inherited}.
\begin{theorem}\label{q198:family}
There is one fixed integer $N_0\ge1$ such that
\[
 \forall N\ge N_0\ \forall x\in\Q,\qquad
 x\in\Z\quad\Longleftrightarrow\quad
 \exists\mathbf y\in\Q^{197},\ P_N(x,\mathbf y)=0.
\]
Each $P_N$ has degree four, content one and 2945 nonzero coordinate
monomials, and all 198 coordinates occur. Its coefficient height is
\[
                         H(P_N)=\max(H_0,A_N^2).
\]
The second branch attains the maximum for $N\ge53$; this coefficient
crossover is not asserted to be a sufficient defining threshold.
The argument here does not compute $N_0$.
\end{theorem}

The new arguments are the additive contact circuit and recovered
point differences of \cite{q198-Brox212,q198-Brox204}, the
constant-positive-weight trace theorem and half-generator reduction
of \cite{q198-BroxWeighted}, and Conant's affine fourth label and
inverse Vieta scalar chart \cite{q198-Conant198}. They are proved
below on their actual domains. The supplying arithmetic foundation
remains that of this manuscript.

\subsection{Exact inherited interfaces}\label{q198:inherited}
Write $K=\Q(\rho)$, $\rho^2=2$, $\sigma(\rho)=-\rho$,
$d=3-\rho$, $D=d^2$, and
\[
 E:y^2=x^3-Dx,\qquad R_0=(1+\rho,2\rho),\qquad P_0=2R_0.
\]
The following results are used with their stated hypotheses.
\begin{enumerate}
\item Proposition~\ref{prop:generator}, transported to this model,
gives $2E(K)=\Z P_0$ with $P_0$ non-torsion and $E(K)_{\rm tors}=E[2](K)$.
Thus every $Q\in E(K)$ is uniquely $[m]R_0+T_0$ with
$m\in\Z$ and $T_0\in E[2](K)$. Put $\nu(Q)=m/2$.
Lemma~\ref{v15:aggregate:lem:independent} gives independence modulo
torsion over $\Z[i]$ for the two pointed conjugate images of $P_0$.
Doubling a relation gives the same independence for the images of $R_0$.
\item Sections~\ref{v222:reference} and~\ref{q214:raw224-source}
recover a finite reference point $H$, with $f(H)=1/x(H)\ne0$, from
the retained polynomial rows before any scalar quotient is interpreted.
Section~\ref{v15:base248:sec:scalar} gives the forced-149 exclusion
of the finite exceptional reference points. Their reciprocal-norm
values are 149-units. It follows, once the scalar valuation test has
been applied, that $H,H-P_0$ are non-torsion and
$S=y(H)/x(H)$ and $T_H=y(H-P_0)/x(H-P_0)$ are defined and nonzero.
These are the raw224-side maps, before the later paired-root and
common-quotient reductions.
\item Section~\ref{v222:contacts} proves for
$\beta=51X^2+2504$ and
$r_s=\beta(\beta-1)(\beta-2401)(32\beta+343)$ that the physical
normalizer is $R=A_Nr_s/(1+149^2r_s)>0$.
Its valuation at 149 is $2N+1$ or $2N-1$.
Every prime in the scalar odd-contact radical $M_5(\beta)$ gives a
positive odd valuation of $R$; these primes exceed 76176.
For integer $X$, $R(R+9792)$ has square reduced denominator and is
a square in $\Q_2,\Q_3,\Q_{17}$.
\item Lemma~\ref{v15x:scalar-local} applies to the physical expression
\[
 W_s=Z_0^4/R^{10}+Z_1^2/R^5+F^2/R^4,\qquad
 b_s=R(R+9792)W_s^2,\qquad F=\Nm(f(H))\ne0.
\]
It allows either component of $Z$ to vanish and does not presuppose
a point interpretation of $Z$. The scalar height supplier
\ref{v15:height:thm:contactheight} gives
$h(\beta)\le C_3M_5(\beta)^3(1+\log M_5(\beta))$,
where $C_3=2^{330770}$.
\item Theorem~\ref{v15:logs:thm:uniform-logarithm-tail} and
Corollary~\ref{v15:logs:cor:uniform-logarithm-ratio} give the
formal-logarithm isometry and relative error of order four at the
consumed good unramified odd primes. The convention $z(O)=0$ is
allowed. Section~\ref{q214:factor-source} proves the valuation
separation and cleared-trace extraction for the five factors used
below, including negative test valuations, zero tests, and split
or inert places. Its integrality step precedes the use of relative
precision as an absolute congruence.
\item Put
\[
 \begin{gathered}
 \Theta_0=\{0,-1/2,-1,1/2399,-32/407\},\qquad
 \Theta=\Theta_0\cup(2-\Theta_0),\\
 \Sigma=\{p:p\le76176\}\cup\{2029147\}.
 \end{gathered}
\]
For $t\in\Q\setminus\Theta$, let $M(t)$ be the squarefree product
of primes outside $\Sigma$ at which some $t-\theta$ has positive
odd valuation. Section~\ref{q214:height-supplier}, with the
operative product argument of Section~\ref{q214:height-averaging},
gives a fixed $H_2>0$ such that
\[
                         h(t)\le H_2M(t)^{7/6}.
\]
Its full domain includes $t=1$, where $M(t)=1$.
\item Theorem~\ref{thm:v5-simultaneous-primes} and
Lemmas~\ref{v15:dependencies:lem:arbitrary-avoidance}
and~\ref{v15:dependencies:lem:rational-contacts} permit any fixed
finite distinct rational coefficient list after square-denominator
clearing, prescribed finite divisibility, finite avoidance, and
arbitrarily large fixed prime and slope bounds. The reduction-order
pair is given by Lemma~\ref{lem:v5-two-reductions}.
Corollary~\ref{v15:dependencies:cor:compatible-progressions} and
Lemma~\ref{v15:dependencies:lem:real-progression} supply the compatible
reference progression and its real recurrence. All these data must
be fixed before the analytic parameter tends to infinity.
\item Lemma~\ref{v5:preparation-zero-bridge}, in the form of
Section~\ref{v15:dependencies:sec:poles}, excludes negative odd
valuations outside $\{2,3,17\}$ for a nonzero input with a pole
certificate. Conversely it supplies a certificate for a rational
$b$ with square reduced denominator and
\[
 0<b<17,\quad b\in\Q_2^{\times2}\cap\Q_3^{\times2}
 \cap\Q_{17}^{\times2},\quad
 v_2(b)\ge3,\quad v_3(b)\ge1,\quad v_{17}(b)\ge2.
\]
Its restricted-converse and descent hypotheses are the ones proved
in the retained foundation. Lemma~\ref{v15:dependencies:lem:conic-complete}
and the retained support charts include zero auxiliary witnesses.
Certificates for different pole inputs have independent coordinates.
\end{enumerate}
Here $h$ on a rational number denotes absolute logarithmic Weil
height. A letter $h$ used as the reference index will be identified
explicitly. These inherited interfaces are not weakened by passing
to the 198 table.

\section{Table transport, contact circuits and point recovery}
\label{q198:table-transport}
% Inventory B09, B01, B03, N01, N02, N04, B04, N03, N06.

\subsection{Transport from the specified predecessor}\label{q198:transport}
The source chain is the raw224-based 214 table, the additive 212 table,
the reciprocal 204 table, Conant's affine-200 intermediate, and the
present 198 table. It does not use the later shifted-200 construction
or an order-four-shift scalar function.

In the 214-to-212 step the two contact cache coordinates and their
defining rows are deleted. The four contact equations for each
detector are replaced by the circuit in Lemma~\ref{q198:contact-circuit}.
The second detector argument becomes unscaled, and the aggregate
relates the two arguments additively. The counts become 212 coordinates
and 186 rows. The contact weight is unchanged physically, as proved
in that lemma. The new additive aggregate has its own arithmetic
proof; it is not an all-witness equivalence to the old aggregate.

In the 212-to-204 step the first and third seven-field root blocks
are replaced by the five-field pairs of Lemma~\ref{q198:reciprocal-pairs}.
Each replacement keeps nine field equations after the ratios are
included. In the 212 numbering the deleted blocks are 155--168 and
183--196, with rational rows 104--121 and 140--157. The ratio
multipliers are respectively 1 and 2. The old coordinates satisfy
$x_{i,\rm old}=R^{e_i}x_{i,\rm new}$, where
\[
 e_i=1\quad(39\le i\le46\text{ or }48\le i\le51),\qquad
 e_i=-1\quad(205\le i\le208),
\]
and the other exponents vanish. Collecting powers of $R$, dividing
a row by its least such power, and replacing even powers using the
retained coordinate $R_2=R^2$ gives the 204 table. The two ratio
square-output pairs use $u_0+2\rho u_1$, whereas ordinary field pairs
use $u_0+\rho u_1$. These conventions are retained in every
component multiplication. Nonvanishing of $R$ is established below
before these substitutions are inverted.

The fourth-block replacement deletes fourteen rational coordinates
and appends ten, with eighteen rows on either side. The final scalar
replacement deletes sixteen rational coordinates and appends fourteen.
It removes 23 rational rows, adds the twenty component equations of
the inverse Vieta chart and its ratios, and adds the new aggregate
row. Thus
\[
 204-14+10-16+14=198,\qquad
 186-18+18-23+20+1=184.
\]
The complete retained-coordinate weights and row division powers
are in the pinned table, rather than inferred from these counts.
The scalar equation $R_2=R^2$ is retained literally, not simplified
away while other rows are divided. The input has weight zero.
The finite identities in Section~\ref{q198:finite-certificate}
verify every retained and changed row. The arguments below give
their arithmetic interpretation.

\begin{lemma}[Compressed contact circuit]\label{q198:contact-circuit}
Let
\[
 P_5(r)=\prod_{\theta\in\Theta_0}(r-(1-\theta)^2)
 =r^5-e_1r^4+e_2r^3-e_3r^2+e_4r-e_5.
\]
With $qs=1$, the four quadratic equations
\[
 r=(t-1)^2,\quad u=r^2,\quad
 v=ru-e_1(1+s)u+e_2(1+s)r,\quad
 uv=(1+s)(e_3u-e_4r+e_5)
\]
are equivalent to $r=(t-1)^2$ and $(1+q)P_5(r)=r^5$.
The auxiliary outputs are unique, $P_5(r)\ne0$, and $t=1$ is
included, with $q=s=-1$ and $u=v=0$.
\end{lemma}
\begin{proof}
After substituting $u=r^2$ and the displayed expression for $v$,
the last residual is $(1+s)P_5(r)-sr^5$. Multiplication by the
nonzero $q$ proves the equivalence, and all substitutions reverse.
If $P_5(r)=0$, the weight equation gives $r=0$, contradicting
$P_5(0)=-e_5\ne0$. At $r=0$ the same equation gives $s=-1$.
The five squared distances are
$1,9/4,4,5750404/5755201,192721/165649$, so $e_5\ne0$.
Hence $t$ avoids all ten elements of $\Theta$.

When $t\ne1$, write
$\mathcal R=P_5((t-1)^2)/(t-1)^{10}$. The weight is exactly
$q=\mathcal R^{-1}-1$. Thus the circuit preserves the physical
contact function, not merely its vanishing set. At $t=1$ no division
by $t-1$ is made; its radical is one and there is no consumed
contact prime. For comparison with the old stored circuit, its first
inverse parameter is $1/t$, with contact output
$t^{-10}P_5((t-1)^2)$. For a previously scaled second argument,
put $w=R^2$ and recover its stored outputs as
$wt,wr,wr^2,wr^3,w^2P_5(r)$. Substitution gives the old contact
identities. This proves the contact portion of the source transport.
\end{proof}

\begin{lemma}[Universal reciprocal pair]\label{q198:reciprocal-pairs}
For $Q\ne(0,0)$ put $f(Q)=1/x(Q)$ and
$\eta(Q)=-y(Q)/x(Q)^2$, with $f(O)=\eta(O)=0$.
The five field quadrics
\[
 F=f^2,\quad G=g^2,\quad \eta^2=f-DfF,\quad
 s(f-g)=\eta+\kappa,\quad
 s(\eta-\kappa)=1-D(F+fg+G)
\]
describe exactly the ordered pairs $U,V\in E(K)\setminus\{(0,0)\}$
with $U\ne V$. Their finite difference is
\[
 Q=U-V=(x_Q,y_Q),\quad
 x_Q=s^2+D(f+g),\quad y_Q=D\eta-s(Df+x_Q).
\]
This includes identity numerator points, tangencies and finite
two-torsion differences.
\end{lemma}
\begin{proof}
The reciprocal equation is $\eta^2=f-Df^3$. Its map
$(f,\eta)\mapsto(-Df,-D\eta)$ is translation by $(0,0)$ on the
original curve. Put
\[
 e_1=s(f-g)-\eta-\kappa,\quad
 e_2=s(\eta-\kappa)-1+D(f^2+fg+g^2),\quad
 e_3=\eta^2-f+Df^3.
\]
Polynomial expansion gives
\[
 \kappa^2-g+Dg^3=e_3-(f-g)e_2+(\eta-\kappa)e_1,
\]
\[
 y_Q^2-x_Q^3+Dx_Q
 =D^2e_3-D(Df+x_Q)(e_2-se_1).
\]
Both reciprocal points and the asserted difference therefore lie on
their curves. If $f\ne g$, $s$ is the chord slope for the translated
first point and the negative translated second point, which gives
the displayed difference. If $f=g$, then $\kappa=-\eta$.
The case $\eta=0$ would require both $f=0$ or $Df^2=1$, and
$1-3Df^2=0$, which is impossible. Thus
$2s\eta=1-3Df^2$ is exactly the tangent formula. Conversely the
chord or tangent construction supplies these equations and its
unique slope. The reciprocal origin represents $O$; no identity
case was divided out. The impossible case $U=V$ is precisely the
vertical chord case excluded by these equations.
\end{proof}

\begin{lemma}[Nonzero scalar scale]\label{q198:nonzero-scale}
Every rational solution of the 198 rows has $R>100$ and $R_2=R^2$.
Consequently all stored Laurent rescalings by $R$ are reversible.
\end{lemma}
\begin{proof}
Let $J_s$ be the stored scalar contact product, delivered coordinate
104. The retained input and product rows, without any division, give
\[
 RJ_s=(51X^2+103)(51X^2+2503)(51X^2+2504)
                   (1632X^2+80471)=r_s>0.
\]
Thus $R\ne0$. The retained forced-prime row is
$1+22201RJ_s-A_NJ_s=0$, whence
$R=A_Nr_s/(1+22201r_s)$. Since $r_s\ge1$ and $N\ge1$,
$R\ge149^3/(1+149^2)>100$. Its square row gives $R_2=R^2$.
\end{proof}

\begin{lemma}[Inverse Vieta chart]\label{q198:inverse-vieta}
Use field coordinates $U,V,W,M,r,p,r_2$, where $p$ is a coordinate,
not the fixed point $P_0$. Impose
\begin{gather}
 UV=R_2p,\quad pW=r,\quad
 r(U+V+W+2M)+DR_2=0,\quad R_2p+W(U+V)=M^2,
 \label{q198:vieta-core}\\
 r_2=r^2,\qquad F_Ur_2=U^2,\qquad F_Vr_2=V^2.
 \label{q198:vieta-outputs}
\end{gather}
These equations give $r,U,V,W,p\ne0$ and finite points
\[
 U_2=\left(\frac{VW}{R^2},\frac{(M+U)VW}{R^3}\right),\qquad
 V_2=\left(\frac{UW}{R^2},-\frac{(M+V)UW}{R^3}\right),
\]
\[
                 Q_2=\left(p,-\frac{(M+W)p}{R}\right),
                 \qquad Q_2=U_2-V_2.
\]
Their abscissas are nonzero and $F_U=f(U_2)^2$, $F_V=f(V_2)^2$.
Conversely every such triple with finite nonzero abscissas supplies
the chart, including its tangent cases.
\end{lemma}
\begin{proof}
Since $DR^2\ne0$, the third equation forces $r\ne0$; the first
two then force $p,W,U,V\ne0$. The three proposed abscissas have
sum $(M/R)^2$, pairwise sum $-(D+2Mr/R^2)$ and product $(r/R)^2$.
For example the pairwise sum is
$r(U+V+W)/R^2=-D-2Mr/R^2$. They are therefore exactly the roots,
with multiplicity, of
\[
 x^3-(M/R)^2x^2-(D+2Mr/R^2)x-(r/R)^2.
\]
This is the intersection polynomial of $E$ with the line
$y=(M/R)x+r/R$. The three points on that line are
$U_2,-V_2,-Q_2$; the chord-and-tangent law gives their sum zero.
Repeated roots cause no division by an abscissa difference.
Moreover $x(U_2)=r/U$ and $x(V_2)=r/V$, proving the output identities.

Conversely the line through $U_2,-V_2,-Q_2$ is nonvertical because
$Q_2$ is finite. Its intercept is nonzero: otherwise $(0,0)$ would
be one of its three intersections, contrary to the three nonzero
abscissas. Write the line $y=s_0x+b_0$ and set
\[
 U=Rb_0/x(U_2),\quad V=Rb_0/x(V_2),\quad W=Rb_0/x(Q_2),
 \quad M=Rs_0,\quad r=Rb_0,\quad p=x(Q_2),
\]
with $r_2=r^2$ and the actual reciprocal-square outputs. Vieta's
identities give every displayed equation. They hold with intersection
multiplicity in the tangent case as well.
\end{proof}

\subsection{Recovery order and ratio domains}\label{q198:recovery-order}
Lemma~\ref{q198:nonzero-scale} first permits the inverse rescalings.
The retained reference identities then give finite $H$ and
$F=\Nm(f(H))\ne0$. At this stage the two rational components of
$Z=RZ_{\rm st}$ are just physical scalar coordinates.
The last term of $W_s$ makes $W_s>0$ and $b_s>0$ regardless of
whether those components vanish. Applying
Lemma~\ref{v15x:scalar-local} at 149 gives
$v_{149}(F)>7v_{149}(R)/4>0$. The finite exceptional reference
norms from Section~\ref{v15:base248:sec:scalar} are
\[
 7,\quad 41^2/2^6,\quad 2^6 7^2/41^2,\quad
 3^4 7^3/17^2,\quad 17^2/(3^4 7).
\]
They are 149-units. The source reference recovery therefore excludes
the torsion and shifted-torsion exceptions and gives $S,T_H\ne0$.
Only now do we interpret the ratio equations below. No conclusion
about $Q_2$ or the physical meaning of $Z$ was used to prove this
reference-domain assertion.

\begin{lemma}[Ratios for reciprocal pairs]\label{q198:pair-ratios}
For a reciprocal pair, a nonzero reference parameter $S$, and a
nonzero rational multiplier $k$, the equations
\[
 \eta L+kfS=0,\qquad (1-DF)L+k\eta S=0
\]
give $L=kz(U)/z(H)$, where $z=-x/y$, on the finite non-torsion
numerator domain. At $U=O$ they give $L=0$. They exclude every
nonidentity two-torsion numerator. The opposite equations with
$(g,\kappa,G,T_H)$ give the identical conclusion for $V,H-P_0$.
Replacing L everywhere by L-1 gives the affine ratio with offset 1.
\end{lemma}
\begin{proof}
If $\eta\ne0$, the first equation gives $L=-kfS/\eta$,
which equals $kz(U)/z(H)$ by the definitions; the reciprocal curve
equation gives the second equation. At the reciprocal origin the
second equation is $L=0$. If $\eta=0$ and $f\ne0$, the first
equation contradicts $kS\ne0$. This also proves the converse for
every allowed numerator and proves the affine assertion by substitution.
\end{proof}

\begin{lemma}[Scalar ratios and non-torsion]\label{q198:scalar-ratios}
Adjoin to the inverse Vieta equations
\[
 L_2(M+U)=2RS,\quad L_2^0(M+V)=-2RT_H,\quad
 Z_{\rm st}(M+W)=-2S.
\]
Then $U_2,V_2,Q_2$ have nonzero ordinates and
\[
 L_2=2z(U_2)/z(H),\quad L_2^0=2z(V_2)/z(H-P_0),\quad
 Z=RZ_{\rm st}=2z(Q_2)/z(H).
\]
In particular $Q_2$ is non-torsion and
$a=2\nu(Q_2)\in\Z\setminus\{0\}$.
\end{lemma}
\begin{proof}
The three right sides are nonzero by Section~\ref{q198:recovery-order}.
Thus $M+U,M+V,M+W$ are nonzero, and the ordinates in
Lemma~\ref{q198:inverse-vieta} are nonzero. Substituting those
ordinates gives the three ratio identities with their displayed
signs and factors of R. All torsion is two-torsion, so $Q_2$ is
non-torsion. Its index is a nonzero half-integer. The stored scalar
quotient remains divided by R; no further normalization coordinate
or inverse guard has been introduced.
\end{proof}

\subsection{The odd fourth label}\label{q198:affine-fourth}
Use reciprocal pairs for $i=1,3,4$. Their physical ratios are
\[
 L_i=k_i z(U_i)/z(H)+e_i,\qquad
 L_i^0=k_i z(V_i)/z(H-P_0)+e_i,
\]
where $(k_1,e_1)=(1,0)$, $(k_3,e_3)=(2,0)$ and
$(k_4,e_4)=(4,1)$. Lemma~\ref{q198:pair-ratios} gives four
quadratic ratio equations for each pair, including the value 1 at
a fourth-block identity numerator. For their finite differences set
\[
 n=\nu(Q_1)\in\tfrac12\Z,\quad a=2\nu(Q_2)\in\Z\setminus\{0\},
 \quad t=2\nu(Q_3)\in\Z,\quad b=4\nu(Q_4)+1\in2\Z+1.
\]
Thus b is nonzero even if $Q_4$ is finite torsion. Reuse the two
Lucas multiplication outputs for $(L_2-1)L_4$ and $(L_2+1)L_3$.
The resulting test is
\begin{equation}\label{q198:lucas-test}
 G=L_3^2-(L_2+1)(L_2-1)L_3L_4+(L_2-1)^2L_4^2-1.
\end{equation}
Indeed if the two outputs are U and V, this expression is
$L_3^2-UV+U^2-1$. Their defining equations and this last expression
are quadratic in the stored coordinates. The affine label costs
no additional coordinate and imposes no false nonzero-index condition
on $Q_4$.

\section{Trace height and local soundness}\label{q198:height-local}
% Inventory B08, N08, B02, B05, N05, N09, N07, B10, N10.

\begin{theorem}[Constant positive weighted trace height]\label{q198:trace-height}
Fix positive rational weights $w_1,\ldots,w_s$. There are constants
$A_*>0,J_*\ge0$, depending only on these weights and the fixed curve
and subgroup, such that every tuple of finite points
$Q_1,\ldots,Q_s\in E(K)$ satisfies
\[
 \max_i|\nu(Q_i)|^2\le A_*h(S_*)+J_*,\qquad
 S_*=\frac14\Tr\sum_i w_ix(Q_i).
\]
The conclusion includes negative or repeated indices, half-indices,
torsion translates and finite points of index zero. The proof does
not make the additive constant effective.
\end{theorem}
\begin{proof}
We prove the constant-weight specialization of
\cite[Theorem 1.1, Lemma 3.1]{q198-BroxWeighted}.
Over a fixed extension, send the curve and its conjugate by pointed
isomorphisms to $E_0:y^2=x^3-x$, using
$(x,y)\mapsto(x/d,y/d^{3/2})$ and its conjugate.
Call the images of $R_0,R_0^\sigma$ respectively $R_*,R_*'$.
They are independent modulo torsion over $\Z[i]$ by the inherited
independence lemma and doubling.

First fix the set of zero indices and their finite two-torsion points.
There are finitely many possibilities. Omit these coordinates from
the abelian ambient variety; they contribute a fixed constant $c$.
If all indices are zero the assertion holds immediately. Otherwise
the remaining tuple in $E_0^{2k}$ has coordinates
$[m_i]R_*+T_i$ and $[m_i]R_*'+T_i'$, where $m_i\ne0$ and the
$T_i,T_i'$ are torsion. Its trace is the value of
\[
 F=\frac d4\sum_iw_ix_i+\frac {d^\sigma}4\sum_iw_ix_i'+c.
\]
We claim no positive-dimensional abelian translate through this tuple,
meeting the domain of F, has constant restriction of F.

Every positive-dimensional abelian subvariety of $E_0^{2k}$ contains
an elliptic subvariety isogenous to $E_0$, by complete reducibility.
An isogeny from $E_0$ into such a subvariety gives a nonzero direction
whose coordinate endomorphisms are $\alpha_i,\alpha_i'\in\Z[i]$.
If two nonconstant coordinates in the first block have a common pole
on the translate of this direction, eliminating the common parameter
gives a torsion relation
$(m_i\alpha_j-m_j\alpha_i)R_*$. Thus
$m_i\alpha_j=m_j\alpha_i$. Since the principal part of the
abscissa at O is $z^{-2}$, the sum of the principal coefficients
at this common pole is
\[
             \frac d{4\alpha_i^2}\sum_jw_j(m_i/m_j)^2\ne0.
\]
It is a nonzero factor times a positive rational sum. Negative and
repeated indices cause no change. A pole cannot be shared across
the two blocks, since that would give a nontrivial relation between
$R_*$ and $R_*'$ modulo torsion. The second block has the same
positive-sum argument. Some coordinate of the direction is nonzero,
so some pole is present and cannot cancel. This proves the claim.

Let $\mathcal V$ be the reduced closure of the graph of F in
$E_0^{2k}\times\mathbb P^1$. It is integral, since the graph over
the nonempty rational-function domain is integral. The actual tuple
belongs to this domain. Fix the finite-rank subgroup generated by
$R_*$ in each first-block coordinate, $R_*'$ in each second-block
coordinate, and the finitely many torsion points. It is independent
of the indices. A positive-dimensional translate contained in the
projective graph fiber through the actual point meets the domain
densely and would make F constant there. The claim excludes it,
including any putative exceptional translate arising at the boundary.

Apply the family height theorem in
\eqref{v15:aggregate:eq:VB}, namely
\cite[Theorem 1.1, printed p.~308]{v15-aggregate-VonBuhren}.
For its embedding hypothesis, embed each $E_0$ by $\mathcal O(3O)$
and take the Segre product. The plane-cubic sequence
\[
 0\longrightarrow\mathcal O_{\mathbb P^2}(n-3)
 \longrightarrow\mathcal O_{\mathbb P^2}(n)
 \longrightarrow\mathcal O_{E_0}(3nO)\longrightarrow0
\]
gives surjective restriction on sections for every $n\ge1$, since
$H^1(\mathbb P^2,\mathcal O(n-3))=0$.
Tensoring and applying Segre verifies the same condition for the
product. The resulting line bundle is ample and symmetric. The
parameter space $\mathbb P^1$ is projective and integral; its height
is $h(S_*)$. The chosen subgroup belongs to the theorem's allowed
set by taking the error point zero. All hypotheses are therefore
satisfied, and its canonical product height is at most a fixed
multiple of $h(S_*)$ plus a fixed constant.

By pointed-isomorphism and Galois invariance this canonical product
height is a fixed positive multiple of
$2\sum_i\widehat h_E(Q_i)$, and
$\widehat h_E(Q_i)=\nu(Q_i)^2\widehat h_E(P_0)$ with
$\widehat h_E(P_0)>0$. It bounds the largest squared index.
Take maxima of the constants over the finitely many frozen-torsion
cases. This proves the assertion on the whole stated domain.
\end{proof}

\begin{corollary}[The actual aggregate]\label{q198:shifted-aggregate}
The new aggregate row gives
\[
 T=7+\frac14\Tr\bigl(x(Q_1)+p+x(Q_3)+x(Q_4)\bigr),\qquad
 t_0=\lambda,\quad t_1=\lambda+T.
\]
Here $p=x(Q_2)$, so $T\ge1$. With
$B=\max(1,|n|,|a|,|t|,|b|)$ there are fixed constants
$A_0>0,J_0\ge0$ such that $B^2\le A_0h(T)+J_0$.
\end{corollary}
\begin{proof}
Each of $x(Q_1),x(Q_3),x(Q_4)$ is the quadratic expression of
Lemma~\ref{q198:reciprocal-pairs}, and p occurs linearly.
Consequently the displayed aggregate is one rational quadratic row.
At either real embedding, $x(x-d)(x+d)\ge0$ gives $x\ge-d$ for
every finite point. Since $d+d^\sigma=6$, the sum of four traces
is at least $-24$, proving $T\ge1$.
Theorem~\ref{q198:trace-height} with all weights one applies to
$S_*=T-7$. For a reduced fraction $u/v$, the transformation
$u/v\mapsto(u-7v)/v$ increases the maximum of the numerator and
denominator by at most a factor eight. Hence
$h(T-7)\le h(T)+\log8$. Finally
\[
 B^2\le32\max_i|\nu(Q_i)|^2+2
\]
by the definitions of the four labels. Absorb these fixed factors
and the height of the translation into $A_0,J_0$. There is no
term involving $h(R)$.
\end{proof}

\begin{lemma}[Additive height transport]\label{q198:additive-height}
If a nonzero integer J satisfies $M(t_j)^k\mid J$ separately
for $j=0,1$, where $k\ge1$, then
\[
                 B^2\le C|J|^{7/(6k)}+J_1
\]
for fixed $C>0,J_1\ge0$ independent of N, the input and witnesses.
No coprimality of the two radicals is required.
\end{lemma}
\begin{proof}
The inherited full-family height inequality and
$h(t_1-t_0)\le h(t_1)+h(t_0)+\log2$ give
\[
 B^2\le A_0H_2\bigl(M(t_0)^{7/6}+M(t_1)^{7/6}\bigr)
                   +J_0+A_0\log2.
\]
Each divisibility individually gives $M(t_j)^k\le|J|$.
Take $C=2A_0H_2$ and $J_1=J_0+A_0\log2$.
\end{proof}

\begin{lemma}[Scalar extraction]\label{q198:scalar-extraction}
Every solution of the 198 rows satisfies, with $c=|a|$,
\[
 149^{L_N}M_5(\beta)^3\mid c,\qquad
 c\ge149^{L_N},\qquad
 h(\beta)\le C_3 149^{-L_N}c(1+\log c/3),
\]
where $L_N=\lfloor9(2N-1)/4\rfloor+1$ and $C_3=2^{330770}$.
\end{lemma}
\begin{proof}
Let $p=149$ or a consumed scalar contact prime, and write
$\mu=v_p(R)>0$ odd. The reciprocal-abscissa valuation of every
finite point is even: if $v(x)>0$, then $2v(y)=v(x)$; if
$v(x)<0$, then $2v(y)=3v(x)$; zero valuation is already even.
The norm valuation $v_p(F)$ is therefore even at both split and
inert primes. Since 9792 is a p-adic unit, the pole obstruction
on $b_s$ gives $v_p(W_s)>-\mu/2$. The three finite summand
valuations in $R^4W_s$ belong respectively to
$2+4\Z$, $1+2\Z$ and $4\Z$. They cannot cancel; zero summands
are omitted. Each exceeds $7\mu/2$, giving
\[
 v_p(Z_0)>19\mu/8,\quad v_p(Z_1)>9\mu/4,\quad v_p(F)>7\mu/4.
\]
These are the hypotheses and conclusion of the inherited scalar
lemma, now applied to the actual doubled quotient of
Lemma~\ref{q198:scalar-ratios}.

At some place over p, $f(H)$ has positive valuation, so H is formal.
Both components of Z have valuation greater than $9\mu/4$, hence
the same holds for Z at that place. The identity
$Z=2z(Q_2)/z(H)$ shows $v(z(Q_2))>0$.
Thus $Q_2$ is formal or reduces to $(0,0)$; in the latter case
translate by that two-torsion point. The translation is
$x'=-D/x$, $y'=Dy/x^2$, and $z'=-z$.
The translated point is formal and has the same index.

Put $h=\nu(H)\in\tfrac12\Z$. The reference exclusions give
$h\ne0$, and h is p-integral. Doubling a suitable group relation
kills the two-torsion and gives the exact formal logarithmic ratio
$2\nu(Q_2)/h=a/h$. In the translated case only the harmless
sign of z changes. The logarithm is an isometry, so
$v_p(a/h)=v(Z)>9\mu/4$, and therefore
$v_p(a)>9\mu/4$. At 149 this gives $v_{149}(a)\ge L_N$;
each prime in $M_5(\beta)$ contributes at least three. These primes
are distinct from 149. The integer a is nonzero by
Lemma~\ref{q198:scalar-ratios}, proving the divisibility and lower bound.
Now $M_5(\beta)^3\le149^{-L_N}c$ and
$\log M_5(\beta)\le(\log c)/3$. Substitute into the inherited
scalar height inequality to obtain the last assertion.
\end{proof}

\subsection{The five factors and half-index comparison}
\label{q198:half-index-comparison}
Put $L_i'=\sigma L_i^0$, $C_i=L_i-L_i'$, $I=L_1-\beta$,
$D_H=f(H)$, $A_i=f(U_i)$, $B_i=\sigma f(V_i)$ and
$B_0=\sigma f(H-P_0)$. At each detector the retained five factors are
\begin{align*}
 &1+q^5C_2^2+q^2D_H^2+q^4B_1^4,\\
 &1+q^5G^2+q^2A_3^2+q^4B_0^4,\\
 &1+q^5I^4+q^3A_2^2+q^2A_4^2+q^4B_3^4,\\
 &1+q^5C_3^2+q^2C_1^4+q^4B_4^4,\\
 &1+q^5C_4^2+q^2A_1^2+q^4B_2^4.
\end{align*}
Here G is the new expression~\eqref{q198:lucas-test}.
At a consumed prime p, precisely one contact has positive odd
valuation $\mu$ and $t_j-1$ is a unit. The contact weight identity
therefore gives $v_p(q)=-\mu$.
All finite reciprocal valuations are even, including those from
the inverse Vieta pair, and identity numerators have reciprocal zero.
Thus the valuation classes in the inherited extraction proof are
unchanged. That proof, applied to the actual displayed tests, gives
at one place over p
\[
 v(C_2,C_3,C_4,G)\ge3,\quad v(C_1,I)\ge1,\quad
 v(D_H,A_i,B_i,B_0)\ge2.
\]
No positivity of test valuations was assumed in obtaining them.

We give the changed formal step explicitly. Write
$h=\nu(H)=m/2$, $u=\nu(U_i)$ and $v=\nu(V_i)$.
The point differences have index $u-v=\nu(Q_i)$.
At the successful place the primary and conjugate opposite points
are formal, with parameter valuation at least one. The identities
\[
 [2m]U_i=[4u]H,\qquad
 [2(m-2)]V_i=[4v](H-P_0)
\]
kill all two-torsion. They give exact logarithmic ratios u/h and
v/(h-1), also for the conjugate curve. At least one of h,h-1 is a
p-adic unit, because p is odd and they differ by one. One exact
ratio is therefore integral. Its relative approximation is integral;
the positive comparison $C_i$ makes the other approximate ratio
integral, and its relative unit factor makes the other exact ratio
integral. Subtracting the fixed integral offset does not alter this
argument, and multipliers 1, 2 and 4 are p-adic units.

After this integrality step, the relative order-four error is an
absolute error of valuation at least four. Let k and e denote the
multiplier and offset, so that both approximations have exact forms
$ku/h+e$ and $kv/(h-1)+e$. Since h and h-1 are integral,
\[
 h(ku/h+e)-(h-1)(kv/(h-1)+e)=k(u-v)+e.
\]
As $L_i=L_i'+C_i$, the left side of the approximate identity is
$L_i+(h-1)C_i$. It follows that
\[
 L_1\equiv n\pmod p,\qquad L_2\equiv a\pmod{p^3},\qquad
 L_3\equiv t\pmod{p^3},\qquad L_4\equiv b\pmod{p^3}.
\]
Identity numerator points are covered by the second ratio equation
and $z(O)=0$. A finite torsion difference is allowed. The same
reasoning applies at inert places; incompatible formal requirements
exclude a solution rather than create an extra branch.

\begin{lemma}[Integer residuals]\label{q198:residual-divisibility}
Put
\[
 \Delta=t^2-(a+1)(a-1)tb+(a-1)^2b^2-1\in\Z.
\]
If $\beta-n\ne0$, let $J_X$ be its positive absolute reduced
numerator. For both j one has
\[
 M(t_j)^3\mid\Delta,\qquad M(t_j)\mid J_X\quad(J_X\ne0).
\]
Furthermore $|\Delta|\le14B^4$ and $J_X\le4B e^{h(\beta)}$.
\end{lemma}
\begin{proof}
The formal congruences and the order-three test G give
$\Delta\equiv0\pmod{p^3}$ at every consumed prime. This is an
integer congruence because a,t,b are integers. The input test gives
$\beta-n\equiv0\pmod p$, hence p divides its reduced numerator.
Multiplying these separate prime divisibilities proves the assertions.
For $B\ge1$, the triangle inequality bounds $|\Delta|$ by
$B^2+(B^2+1)B^2+(B+1)^2B^2+1\le14B^4$.
Write $n=m/2$, $|m|\le2B$, and $\beta=u/v$ in lowest terms.
The unreduced numerator $2u-mv$ has absolute value at most
$4B\max(|u|,v)=4Be^{h(\beta)}$; reduction can only decrease it.
\end{proof}

\begin{lemma}[Affine Lucas growth and converse]\label{q198:affine-lucas}
If $\Delta=0$, b is nonzero, and $c=|a|$ is sufficiently large,
then
\[
                         \log B\ge(c-6)\log(c-3).
\]
Conversely, for every sufficiently large positive $a\in6\Z$, the
recurrence $U_0=0$, $U_1=1$, $U_{j+1}=(a+1)U_j-U_{j-1}$ gives
integers
\[
 j=a-1,\qquad b=U_j/(a-1),\qquad t=U_{j+1}
\]
with b positive and odd, $b>1$, and $\Delta=0$.
\end{lemma}
\begin{proof}
Put $A=a+1$, $c'=|A|$ and $w=\operatorname{sgn}(A)(A-2)b$.
For large c, $c'>2$, $w\ne0$, and the equation becomes
$t^2-c'tw+w^2=1$. We recall the elementary norm-one argument used
in Section~\ref{q214:uniform-existence}, P12. With
$\epsilon=(c'+\sqrt{c'^2-4})/2$, every positive norm-one unit in
$\Z[\epsilon]$ greater than one is a power of $\epsilon$.
Indeed a unit strictly between 1 and $\epsilon$ would have
$u-u^{-1}=b_0\sqrt{c'^2-4}$ for a nonzero integer $b_0$, which is
impossible in that interval; division by an appropriate power of
$\epsilon$ proves the assertion. Apply this to $t-w\epsilon$,
allowing sign and inversion. The recurrence identities give
$|w|=U_j(c')$ and $|t|\in\{U_{j-1}(c'),U_{j+1}(c')\}$ for $j>0$.

Modulo $c'-2$, $U_j(c')\equiv j$; modulo $c'+2$,
$U_j(c')\equiv(-1)^{j-1}j$. If A is positive use the former
modulus $|A-2|$; if A is negative use the latter. Since this
modulus divides w, $j\ge c'-2$ in either case.
Induction in the recurrence gives
$U_j(c')\ge(c'-1)^{j-1}$ for $j\ge1$.
Thus
\[
 \log B\ge(j-2)\log(c'-1)
 \ge(c'-5)\log(c'-2)\ge(c-6)\log(c-3)
\]
for all sufficiently large c, using $c'\ge c-1$.

For the converse A=a+1 is odd and j=a-1 is odd and congruent
to 2 modulo 3. The recurrence modulo a-1 gives $U_j\equiv j=0$,
so b is an integer. Modulo two its terms repeat $0,1,1$; therefore
$U_j$ and b are odd. Positivity and $b>1$ hold for large a.
Finally $U_{j+1}^2-AU_jU_{j+1}+U_j^2=1$, by induction, and
$U_j=(a-1)b$ gives the required equation.
\end{proof}

\begin{lemma}[Recovery of the ordinary input]\label{q198:integer-input}
If $\beta=51X^2+2504\in\tfrac12\Z$ for $X\in\Q$, then X is an
integer.
\end{lemma}
\begin{proof}
For $X=u/v$ in lowest terms, $v>0$, integrality of
$2\beta=102u^2/v^2+5008$ implies $v^2\mid102u^2$.
Coprimality gives $v^2\mid102$. Since 102 is squarefree, $v=1$.
\end{proof}

\begin{proposition}[One soundness threshold]\label{q198:soundness}
There is a fixed $N_S$ such that every rational zero of the 198
table for every $N\ge N_S$ has integer input.
\end{proposition}
\begin{proof}
If $\Delta\ne0$, Lemmas~\ref{q198:additive-height}
and~\ref{q198:residual-divisibility} give
\[
 B^2\le C14^{7/18}B^{14/9}+J_1.
\]
Since $14/9<2$, this bounds B by a fixed constant. Choose a fixed
threshold beyond which $149^{L_N}$ exceeds that bound. The scalar
lower bound then forces $\Delta=0$ for every solution.

If $J_X\ne0$, the same height lemma with k=1, once $B^2\ge2J_1$,
gives
\[
 \log B\le\frac75h(\beta)+C_2,\qquad
 C_2=\frac65\log(2C)+\frac75\log4.
\]
Increase the threshold so $(7/5)C_3 149^{-L_N}\le1$ and c is
large enough for Lemma~\ref{q198:affine-lucas}. Its upper bound is
then at most $c+(c/3)\log c+C_2$, whereas the lower bound is
$(c-6)\log(c-3)$. Their difference tends to positive infinity:
after division by c it is $(2/3)\log c-1+o(1)$.
A further fixed lower bound on c, and hence on N, excludes this
case for every input. Thus $\beta=n\in\tfrac12\Z$, and
Lemma~\ref{q198:integer-input} applies. The constants were fixed
before N, X or any witness. Taking their maximum gives one $N_S$.
\end{proof}

\section{Half-generator reduction and fresh completeness}
\label{q198:completeness-section}
% Inventory B06, B07, N11, B02 completeness, N12.

\begin{lemma}[The half-generator at the selected primes]
\label{q198:half-generator-reduction}
Let $p=A^2+B^2$ be a good prime with
$A\equiv9\pmod{16}$ and $B\equiv0\pmod{16}$.
The reductions of $R_0$ and its conjugate lie in the doubled
subgroups of their respective elliptic groups over $\F_p$.
Thus the quarter of the group order annihilates each reduction.
\end{lemma}
\begin{proof}
This is the reciprocity argument of
\cite[Lemma 5.1]{q198-BroxWeighted}.
We have $AB\ne0$, $\gcd(A,B)=1$ and $p\equiv1\pmod8$.
In $\F_p$ set $i=A/B$, so $i^2=-1$.
Quadratic reciprocity gives $(B/p)=1$: its sign and 2-part are
squares, and for every odd prime $q\mid B$ one has
$(p/q)=(A^2/q)=1$, with $p\equiv1\pmod4$.
Using Jacobi symbols for a composite denominator,
\[
 \left(\frac{A+B}{p}\right)
 =\left(\frac p{|A+B|}\right)
 =\left(\frac2{|A+B|}\right)=1.
\]
Indeed $p\equiv2A^2\pmod{|A+B|}$, while $A+B\equiv1\pmod8$;
absolute value does not change the symbol of 2. Consequently
$2+2i=2(A+B)/B$ is a square. Its product with $2-2i$ is 8,
also a square, so both are squares. Choose roots u and v.
Then $uv=\pm2\sqrt2$ and
$(u\pm v)^2/4=1\pm\sqrt2$, with the signs possibly exchanged.
Thus both $1+\rho$ and $1-\rho$ are squares for either reduction
of $\rho$.

The differences of the abscissa of $R_0$ from the three roots
$0,d,-d$ are
\[
              1+\rho,\qquad 2(\rho-1)=2/(1+\rho),\qquad4.
\]
They are nonzero squares. The full-two-torsion halving criterion
therefore makes $R_0$ a double. Explicitly, choose square roots
$r_i$ of these differences with $r_1r_2r_3=-y(R_0)$;
the abscissa of a half is
$x(R_0)+r_1r_2+r_1r_3+r_2r_3$, with ordinate
$-y(R_0)-(r_1+r_2+r_3)(r_1r_2+r_1r_3+r_2r_3)$.
Substitution and the duplication formula give twice this point equal
to $R_0$. The same calculation works for the conjugate.

A finite elliptic group with full rational two-torsion is
$\Z/a\Z\times\Z/b\Z$, where $a\mid b$ and a,b are even.
The doubled subgroup has exponent $b/2$, which divides $ab/4$.
Thus the quarter-order annihilates every double, proving the claim.
\end{proof}

\begin{lemma}[The half-generator reference progression]
\label{q198:half-generator-progression}
Fix finitely many integer multiples $Q_i=[m_i]R_0$, a scalar
progression $h\equiv h_*\pmod L$ with $64\mid h_*,L$, and
compact real and finite-place neighborhoods for the required point
functions on that progression. Include the real component-order
divisibility in L. For the primes supplied by the inherited
fixed-data theorem, the same order moduli as in the doubled-point
construction permit
\[
 H=[h]P_0,\qquad U_i=[h]Q_i,\qquad V_i=[h-1]Q_i
\]
to have the required primary and opposite formal reductions,
while preserving the prepared and real conditions. Arbitrarily
large h can be used, in particular $h\ne0,1$.
\end{lemma}
\begin{proof}
The two conjugate reduction orders at a selected prime are the
unordered pair $d_-(\pi),d_+(\pi)$. The prime theorem imposes
$\operatorname{Re}\pi\equiv9\pmod{16}$ and
$\operatorname{Im}\pi\equiv0\pmod{16}$.
Lemma~\ref{q198:half-generator-reduction} shows that their quarters
kill $R_0$, hence every $Q_i$. With
\[
 U=\operatorname{lcm}(d_-(\pi)/4),\qquad
 V=\operatorname{lcm}(d_+(\pi)/4),
\]
the inherited avoidance and order calculations give
$\gcd(U,V)=1$, $\gcd(U,L)=16$ and $\gcd(V,L)=1$.
Therefore
$h\equiv0\pmod U$, $h\equiv1\pmod V$,
$h\equiv h_*\pmod L$ are compatible. They put H and every
$U_i$ in the primary formal group and $H-P_0$ and every $V_i$
in the opposite one. The roles of the two places follow the
unordered order pair.

For the real conditions apply
Lemma~\ref{v15:dependencies:lem:real-progression} to the closure
of $(R_0,R_0^\sigma)$ in the product of the real elliptic groups.
Every displayed function is an integer multiplication or a fixed
translate of it: $H=[2h]R_0$, $U_i=[hm_i]R_0$ and
$V_i=[hm_i-m_i]R_0$. No multiplication-by-one-half rational map
is invoked. The compatible progression is dense in the prescribed
component. Its arbitrarily large positive indices meet the selected
nonempty real neighborhoods. Finite-place conditions were included
in L and are preserved. Removing h=0,1 deletes only two indices.
\end{proof}

\begin{proposition}[Uniform scalar preparation]\label{q198:scalar-preparation}
There is a fixed integer $N_C$ such that, for every $N\ge N_C$
and integer input X, one can choose a positive integer $a_0$,
$3\mid a_0$, set $Q_2=[a_0]P_0$, $a=2a_0\in6\Z$, and
choose a reference progression $h\equiv h_*\pmod L$ with
$64\mid h_*,L$ on which the final scalar input satisfies every
hypothesis of the inherited zero-preparation realization theorem.
The index $a_0$ can be arbitrarily large.
\end{proposition}
\begin{proof}
This is the application of P13--P14 in
Section~\ref{q214:uniform-existence} to the actual doubled scalar
quotient, with the finite set and initial divisibility made explicit.
For integer X, $r_s$ is an integer at least one, so
$R\ge A_N/(1+149^2)$ uniformly in X.
Choose the fixed compact real neighborhoods of the inherited
scalar construction. For its undoubled quotient
$Z^{\rm old}=z(Q_2)/z(H)$ and norm F, let
$|Z^{\rm old}_0|,|Z^{\rm old}_1|\le M_Z$ and $|F|\le M_F$
there, with constants independent of X,N.
Since $Z=2Z^{\rm old}$, for $R\ge1$ we have
\[
 0<W_s\le K'_W R^{-4},\qquad
 K'_W=16M_Z^4+4M_Z^2+M_F^2,\qquad
 0<b_s\le9793(K'_W)^2R^{-6}.
\]
Choose a fixed $N_C$ making the last upper bound less than 17
for every larger N. Positivity follows from F nonzero. This
choice precedes all input-dependent moduli and labels.

For the given pair (N,X), put $D_s=1+149^2r_s$.
The four-contact calculation gives
$\gcd(A_Nr_s,D_s)=\gcd(A_Nr_s+9792D_s,D_s)=1$,
so $R(R+9792)$ has square reduced denominator $D_s^2$.
It also gives $v_2(R)\in\{2,3\}$ and $v_3(R)=v_{17}(R)=0$.
Thus $1+9792/R$ is a square at 2,3,17, and so is $R(R+9792)$.
Let $\mathcal S$ contain its odd-valuation numerator support and
\emph{also} 2,3,17. These three primes are included even if they
are absent from that support.

Choose the initial reference $H_*=[h_*]P_0$ sufficiently formal
at every place over $\mathcal S$, with $64\mid h_*$ from the
start, and in the fixed real neighborhoods. This makes F as deep
as required. Next choose $a_0$ sufficiently divisible so that
$[a_0]P_0$ is still deeper at those places, while its real image
lies in the fixed neighborhood of O. Impose $3\mid a_0$ and
any prescribed lower bound on $a_0$. The finite-index compact-group
recurrence used in P14 permits this additional fixed divisibility.
The quotient $z(Q_2)/z(H_*)$ can consequently be made arbitrarily
deep in both rational components. Multiplication by two leaves
odd-prime depths unchanged and increases their dyadic depths.
It follows that each of the three terms of $W_s$ is integral at
the odd-valuation support, and that the scalar input has depths
at least 3,1,2 at 2,3,17. The local squareclass is unchanged,
since $b_s=R(R+9792)W_s^2$.

At the odd-valuation numerator support, integrality of $W_s$
prevents negative odd valuations of $b_s$. Outside that support
the valuation of $R(R+9792)$ is even; multiplication by $W_s^2$
preserves even parity. Thus the reduced denominator of $b_s$ is
square. The real estimate gives $0<b_s<17$, and the local square
and depth conditions have just been verified.

Now fix $a_0$. By continuity, a sufficiently small finite-place
and real neighborhood of $H_*$ preserves all these conditions and
the nonzero reference and scalar-chart denominators. Choose L
with 64 dividing L, including the finite preparation and component
orders, so $h\equiv h_*\pmod L$ preserves the finite conditions;
retain the chosen real neighborhood for the later recurrence.
The initial h* was already divisible by 64, so the later CRT
compatibility is available. These are hypotheses on the scalar
input at the later H; no fixed earlier scalar witnesses are being
carried along this progression.
\end{proof}

\begin{lemma}[The additive coefficient list with coincidences]
\label{q198:additive-contacts}
For fixed rational $T\ge1$, the two detectors use the coefficient
list $\Theta\cup(\Theta-T)$. It is admissible for the inherited
prime-selection theorem after equal coefficients are merged and
every incidence is retained. The merged list has at least ten
distinct coefficients. The construction gives the contact parity,
finite-place preparation and real limits required by both detectors.
\end{lemma}
\begin{proof}
An incidence is a pair $(j,\theta)$, with offsets $a_0=0,a_1=T$
and $\theta\in\Theta$. Associate the coefficient
$\gamma=\theta-a_j$ to it. Within each detector these ten
coefficients are distinct. Across detectors they may coincide;
the bound $T\ge1$ does not exclude that possibility.
Merge them into a finite set and keep the list of incidences for
each member. Coefficient zero occurs, so the selected primes include
the numerator u of the slope.

Choose a positive integer D with $D^2\gamma\in\Z$ for every
coefficient, and apply the inherited theorem to the distinct integers
$c_\gamma=-D^2\gamma$. Prescribed divisibility of w and an enlarged D
give any fixed denominator depths. Set
\[
 v_0=D^2w^2,\qquad\lambda=u/v_0,\qquad
 p_\gamma=u-\gamma v_0=u+c_\gamma w^2.
\]
All coefficients, finite preparation data, forbidden primes and real
neighborhoods are fixed before applying the theorem. Its slope and
prime cutoffs can ensure $p_\gamma>w^2$, avoid all fixed denominators
and coefficient differences, and make all selected primes distinct.
Then each incident contact $t_j-\theta$ has valuation one at
$p_\gamma$, while other contacts in that detector are units.
The same prime and Gaussian generator is used for every incidence
of a merged coefficient. Its reference congruences are imposed once.

For either detector, away from $t_j=1$, its exact normalized contact
function is
\[
 \mathcal R_j=
 \frac{\prod_{\theta\in\Theta}(u+(a_j-\theta)v_0)}
 {v_0^{10}(\lambda+a_j-1)^{10}}.
\]
Every numerator factor is its selected prime, and each occurs once
within that detector. The denominator is a rational square.
Thus all other valuations of $\mathcal R_j$ are even. At a fixed
prepared prime, choose the denominator depth so that $\lambda$
dominates 1 and T p-adically. The identity
\[
 \mathcal R_j=
 \prod_{\theta\in\Theta_0}
 \left(1-\frac{(1-\theta)^2}{(t_j-1)^2}\right)
\]
makes $\mathcal R_j$ arbitrarily close to one and
$q_j=\mathcal R_j^{-1}-1$ arbitrarily small there. The selected
primes avoid the preparation set, so this denominator choice remains
valid. At real places the same expression tends to one from below
as $\lambda\to+\infty$; hence $q_j>0$ tends to zero. Arbitrarily
large fixed slope bounds give these simultaneous conditions.
This argument uses no assertion of twenty distinct coefficients.
\end{proof}

\begin{proposition}[One completeness threshold]\label{q198:completeness}
For every integer X and every $N\ge N_C$, with $N_C$ enlarged
once if necessary by a fixed amount, the 198 rows have rational
witnesses with input X.
\end{proposition}
\begin{proof}
Apply Proposition~\ref{q198:scalar-preparation} and choose $a_0$
arbitrarily large, so $a=2a_0\in6\Z$ is in the range of
Lemma~\ref{q198:affine-lucas}. Choose its positive odd b and t.
Set
\[
 Q_1=[\beta]P_0,\quad Q_2=[a_0]P_0,\quad Q_3=[t]R_0,\quad
 Q_4=[(b-1)/2]R_0.
\]
All four are non-torsion and finite, since $\beta\ge2504$,
$a_0,t>0$ and $b>1$. Their labels are exactly $\beta,a,t,b$,
and their Lucas residual is zero. These points and the rational
$T\ge1$ of Corollary~\ref{q198:shifted-aggregate} are now fixed
before comparison-prime selection.

Choose compact neighborhoods for the finitely many reference and
numerator functions in the product real group, within the scalar
neighborhood and away from their finitely many exceptional points.
They exist because all numerator functions are fixed integer
multiplications and translates, and the initial nonzero integer
reference and its predecessor give non-torsion numerator points.
Include the component orders in L. At the prepared finite places
choose neighborhoods on which those rational functions are bounded.
Enlarge the preparation set to include $\Sigma$, the scalar support,
all fixed denominator and guard primes, and numerator/denominator
primes of T. These data depend on the fixed (N,X) and labels,
which is permitted by the fixed-data prime theorem.

Apply Lemma~\ref{q198:additive-contacts}, choosing sufficiently
large fixed prime and slope cutoffs, and then
Lemma~\ref{q198:half-generator-progression}. Let the final index h
avoid 0 and 1 and meet the retained neighborhoods. Put
\[
 H=[h]P_0,\qquad U_i=[h]Q_i,\qquad V_i=[h-1]Q_i.
\]
The scalar triple consists of non-torsion points with nonzero
abscissas and ordinates, so the inverse Vieta converse applies.
The reciprocal-pair converses apply to the other three pairs.
The reference and all stored ratios have their stated definitions.

At a selected contact prime, the primary formal group contains H
and every $U_i$, and the conjugate opposite formal group contains
$H-P_0$ and every $V_i$. Each reciprocal has valuation at least two.
The exact logarithmic ratios are the chosen labels, with their
multipliers and offset. Since those labels are p-integral, their
relative order-four errors are absolute errors of order at least
four. Therefore all $C_i$, I and G have valuation at least four.
Here $v(q_j)=-1$. The smallest valuations of the nonconstant terms
in the five factors are, according to their type,
\[
 \begin{array}{c|cccccc}
 \text{term}&q^5C_i^2&q^5I^4&q^3A_2^2&q^2A_i^2
                  &q^2C_1^4&q^4B_i^4\\ \hline
 \text{lower bound}&3&11&1&2&14&4.
 \end{array}
\]
The same bound two holds for $q^2D_H^2$, and the G term has
the same bound as $q^5C_i^2$. Thus the factor product $W_j$
is one modulo the successful place. At the opposite place, the
inherited distinct-valuation-class argument, which needs no formalness
assumption there, says $W_j^{-1}$ has residue zero or one.
Consequently $\tau_j=\Tr(W_j^{-1})$ is a unit with residue one or
two at this odd rational prime.

The actual detector pole input is
\[
             b_j=\gamma^2\mathcal R_j/\tau_j^2,\qquad
                         \gamma=408/409.
\]
It has valuation one at its selected primes, and even valuation at
all other primes, by the contact factorization in
Lemma~\ref{q198:additive-contacts}. In particular its reduced
denominator is square. At the prepared primes q is arbitrarily
small uniformly on the fixed neighborhoods, so $W_j$ is close to
one and $\tau_j$ close to two. Choose the closeness so
$\mathcal R_j$ is a square at 2,3,17. The valuations of $\gamma$
there are 3,1,1, giving the exact detector depths 4,2,2 and the
required local squareclasses. At real places q is positive and
small; every factor is positive and tends to one. Thus
$\tau_j>0$ and $b_j\to\gamma^2/4\in(0,17)$.
The fixed slope cutoff was chosen to make $0<b_j<17$ for both
actual detector inputs. It also ensures $q_j\ne0$ and $t_j\ne1$.

The final H still satisfies all scalar conditions of
Proposition~\ref{q198:scalar-preparation}. Now, \emph{after this
final H is fixed}, apply the zero-preparation realization theorem
to the actual scalar input $b_s(H)$ and to the two actual inputs
$b_0,b_1$. Choose their certificates separately. The scalar
conditions were preserved by continuity, but an earlier tuple of
scalar witnesses was not preserved or reused. The complete conic
and support charts allow zero auxiliary witnesses and introduce no
extra restriction.

Fill the contact outputs using Lemma~\ref{q198:contact-circuit},
the point and ratio outputs by their proven converse formulas, and
the retained reference, scalar, factor and trace coordinates by
their physical definitions. Apply the exact inverse stored-coordinate
rescalings, legitimate because $R>0$. Every retained and changed
row then vanishes. All later choices depended on the fixed (N,X)
and labels, but imposed no further lower bound on N. The one fixed
scalar real threshold therefore proves the proposition.
\end{proof}

\section{Exact finite identities and completion}\label{q198:finite-certificate}
% Inventory N13 and the finite part of B09.
The selected unbalanced quadratic table has SHA-256
\begin{center}\small\ttfamily
c55d7eefe41f62fee88b467e3ea8f6be8ba1e66f6a94b30eb74b39f6f6e7efad.
\end{center}
Its expanded polynomial has SHA-256
\begin{center}\small\ttfamily
652b5141aaa0d9f54ef89196ae9e442787335b2fd498e250c2772dc3af98a36c.
\end{center}
These identities fix the table, component conventions, coordinate
weights and primitive row normalizations used here.
The exact reconstruction described in \cite{q198-FiniteReconstruction}
compares every retained row under the complete Laurent weight map,
derives the twenty new inverse Vieta component rows directly from
their displayed field equations, retains $R_2=R^2$ literally,
and reconstructs the aggregate. It independently expands every
quadratic square by an unordered term-pair sum. The resulting
coefficient dictionary is exactly the pinned polynomial dictionary.

\begin{proposition}[Finite coefficient identities]\label{q198:finite-algebra}
The specified 198 table has 184 primitive integral quadratic
residuals at every allowed parameter. Its sum of squares has degree
four, content one, all 198 coordinates present, and 2945 coordinate
monomials. Writing $\mathcal P_0$ for its fixed coefficient table
at A=0, one has
\[
 P_A=\mathcal P_0-2AY_{103}-44402AY_{103}^2Y_{115}
                                      +A^2Y_{103}^2.
\]
The three displayed monomials are absent from $\mathcal P_0$,
which has 2942 nonzero monomials and height $H_0$.
Consequently $H(P_N)=\max(H_0,A_N^2)$.
\end{proposition}
\begin{proof}
We describe the finite certificate and the deductions from it.
An ordinary field pair multiplies by
$(u_0,u_1)(v_0,v_1)=(u_0v_0+2u_1v_1,u_0v_1+u_1v_0)$;
the two specified reciprocal-square slots are converted from their
doubled-second-component convention before this multiplication.
Substitute the recorded Laurent weights, reduce by $\rho^2=2$,
and collect each rational row. Multiply by its denominator and
divide by its integer content with its recorded sign. These
operations give precisely the 184 stored rows. The comparisons
are coefficient identities, not numerical specializations.
The source reference and contact identities in
Lemma~\ref{q198:nonzero-scale} are among these retained identities,
so reversing row division on the zero set does not lose a
zero-scale component.

For a residual $\sum_\alpha c_\alpha X^\alpha$, its square is
computed as the diagonal terms $c_\alpha^2X^{2\alpha}$ and the
off-diagonal terms $2c_\alpha c_\beta X^{\alpha+\beta}$ for
$\alpha<\beta$. Adding all 184 dictionaries gives 2942 constant
profiles, two linear-A profiles with coefficients -2 and -44402,
and one A-squared profile with coefficient 1. Their monomials are
distinct. In delivered one-based numbering the last three are
$x_{104}$, $x_{104}^2x_{116}$ and $x_{104}^2$; since $x_1=X$
and $x_i=Y_{i-1}$, these are exactly the displayed formula.
The largest constant magnitude is $H_0$, attained at $x_{110}^4$.
The coefficient of $X^4$ is exactly one with no parameter term.
These finite integer equalities are recorded by the complete
coefficient comparison; they can also be checked directly from
the supplied dictionaries by the preceding expansion rule.

The $X^4$ coefficient proves degree at least four and integer
content one at every specialization; the quadratic residual degrees
prove degree at most four. Each coordinate appears in the verified
nonzero coefficient dictionary. No allowed positive A can cancel
a coordinate monomial, because every profile contains only one
power of A. This proves the support assertions. The unbalanced
rows are primitive after specialization: their parameter supports
are disjoint, and their coefficient gcds at $A=149^3$ are one.
For larger allowed powers only the same prime 149 can enter a
parameter coefficient. A constant coefficient excludes that prime
whenever necessary, as the gcd at $149^3$ already shows. Any
other prime has unchanged divisibility in the parameter coefficients.
Thus no new common prime divisor appears.

Finally $A\ge149^3>44402$, so the two linear coefficient magnitudes
are at most $A^2$. Both the pure $H_0$ and pure $A^2$ profiles
occur, proving the exact maximum. The integer comparisons
$149^{210}<H_0\le149^{214}$ give crossover N=53 among positive
integers N. This is purely a coefficient statement.
\end{proof}

\begin{proof}[Proof of Theorem~\ref{q198:family}]
For a fixed N, a sum of rational squares vanishes exactly when each
residual vanishes. The table transport and chart identities identify
those rows with the physical system used above, with all its stated
domains. Propositions~\ref{q198:soundness} and~\ref{q198:completeness}
therefore give the claimed input projection for
$N\ge N_0=\max(1,N_S,N_C)$. All constants in this maximum were
fixed independently of the input. Proposition~\ref{q198:finite-algebra}
supplies the coefficient and degree assertions for the same polynomial.
\end{proof}

\begin{remark}\label{q198:bb-interface}
The fixed specialization
$N_{\rm BB}=\operatorname{BB}(2^{1000})$ and the strict bound
$H(P_{N_{\rm BB}})<\operatorname{BB}(2^{2000})$ require a separate
threshold and machine-size proof for this 198 family.
They are not consequences of the crossover N=53.
In this unbalanced presentation, if that specialization is supplied,
the coefficient parameter is exactly
$A=149^{2\operatorname{BB}(2^{1000})+1}$ and every coefficient is
an integer polynomial of degree at most two in A. The family proof
above asserts no numerical sufficient N and no Busy Beaver threshold.
\end{remark}
